\documentclass[10.5pt,a4paper]{article}
\usepackage[margin=1.9cm]{geometry}
\usepackage[T1]{fontenc}
\usepackage{lmodern}
\usepackage{microtype}
\usepackage{parskip}
\usepackage{titlesec}

\titleformat{\section}{\large\bfseries}{\thesection}{0.6em}{}
\titlespacing*{\section}{0pt}{0.9em}{0.3em}

\title{\vspace{-1.5em}\textbf{Matrix-Multiply Accelerator for a RISC-V SoC:}\\
\large Progress Summary}
\author{}
\date{November 2026}

\begin{document}
\maketitle
\vspace{-1.2em}

\section{Work Done}

The project began with a working but very limited accelerator. It used a fixed $4\times 4$
systolic array, handled matrices of at most $4\times 4$, and moved every word over the bus
with its own transaction. Each word cost about 4.7 cycles, and computation was only 4\% of
the total time. The first task was to make the development flow reliable. The simulation
build now runs on Linux, rebuilds correctly after RTL and configuration changes, and keeps
large waveforms and build products out of version control. A memory model with random
ready/valid stalls was added so that handshake bugs would not be hidden by a memory that
always answers at once.

The second task made the array size a build-time choice. The array is now a generated grid
of processing elements, parameterized by \texttt{DIM} (the grid size) and \texttt{MAX\_K}
(the buffer depth along the inner dimension). The same RTL builds arrays from $2\times 2$ to
$16\times 16$, including non-power-of-two sizes such as $5\times 5$. A read-only configuration
register lets the software read the hardware size at run time, so one test program works on
every build.

The third task replaced single-word transfers with AXI4 burst engines. Separate read and
write engines stream whole row segments, with several reads in flight, and respect the
4\,KB boundary rule. This reduced the cycle counts substantially, from 234 to 79 cycles for a
$4\times 4$ product and from 3630 to 877 cycles for a $16\times 16$ product, a speed-up of
about 3 to 4 times. The fourth task added hardware tiling. A single start command now computes
a product of any size by looping over tiles of the result, with strided access so that
sub-matrices can be used, and with A-panel reuse that loads each row of A once for every
tile in its row panel. For a $64\times 64$ product, reuse reduced memory reads by 47\% and
run time by about 36\%.

The design was then brought up on a RealDigital Boolean board, which carries a Xilinx
XC7S50 FPGA. The board top connects the SoC to 128\,KB of block RAM, which holds the
program and all data. A separate software build places the program at \texttt{0x2000}, the
board's reset vector, and writes a memory image that is compiled into the bitstream. The
clock is generated by a PLL that turns the 100\,MHz oscillator into 25\,MHz, and the console
runs over the board's USB-UART at 1\,Mbaud. The Vivado flow is scripted: one script builds
the bitstream and writes timing and utilization reports, and another loads it over JTAG.
With the array at $5\times 5$, the design meets timing at 25\,MHz. At \texttt{MAX\_K=64} it
uses 52\% of the logic cells and 42\% of the flip-flops. The full test suite passes on the
board, including $64\times 64$ products with and without reuse.

\section{Verification Methods}

Every accelerator result is checked against a reference computed by the RISC-V CPU itself.
The reference is a plain triple loop over the same matrices, and the test compares the two
outputs element by element. The test suite covers fixed hand-written inputs, $1\times 1$
products, full and partial tiles, odd sizes that need several K chunks, dot and outer
products, strided sub-matrices, matrices placed just before a 4\,KB boundary, and the error
cases. For the strided case, the test checks that padding is never read and never written.

The suite was run across a sweep of configurations: array sizes of 2, 4, 5, 8 and 16, buffer
depths from 4 to 64, and both single-beat and burst settings, including a configuration with
one outstanding burst and single-beat transfers. Each configuration was also run with random
memory stalls. Verilator linting is clean at several array sizes. For the board, the same
RTL was simulated in a Verilator model of the board top, with a C++ decoder that reads the
UART output bit by bit. The simulated output, including the cycle counts, matches the output
captured from the real board for the $5\times 5$ build with \texttt{MAX\_K=7}, byte for
byte. Synthesis, implementation, and the bitstream were run with Vivado, and the timing and
utilization reports were checked after each build.

\section{Limitations}

The accelerator is limited by memory, not computation. After bursts and tiling, loading
takes about 75 to 85\% of the total cycles, and computation takes 10 to 15\%. The data path is
32 bits wide, and each tile loads $2\cdot DIM\cdot K$ words to perform $DIM^2\cdot K$
multiply-accumulates. So the array can sustain at most about $DIM/2$ multiply-accumulates per
cycle, against a peak of $DIM^2$. For a $4\times 4$ array this caps utilization near 12\%. The
measured $64\times 64$ run reaches roughly 14\% of peak, which is close to that ceiling.

The Boolean board has only 128\,KB of main memory, and the FPGA has about 330\,KB of block
RAM in total. Three $N\times N$ matrices of 32-bit values need $12N^2$ bytes, so the current
test layout runs up to about $N=90$, and $64\times 64$ is the largest size tested on the
board. Products of $256\times 256$ and larger cannot be held on this board at all. A second
limitation of the memory system is that it decodes only the low address bits. An address
beyond 128\,KB wraps around and overwrites the program or the stack, with no error. The test
does not yet check that its buffers fit, so a configuration error could produce wrong results
silently.

Several other limits remain. Software and hardware dimensions are 16 bits wide, so values
above 65\,535 are not supported. The input buffers are built from flip-flops and wide
multiplexers, and their size grows quickly with \texttt{MAX\_K}. At \texttt{MAX\_K=64} the
logic use was already 52\%, so larger buffers will eventually fail to place. The 32-bit
multiplies consume DSP slices, which limits the array to about $DIM=6$ on this device. The
accelerator polls for completion instead of raising an interrupt, and it does not abort a job
cleanly on a bus error.

\section{Pending Work}

The next phase is to remove the memory limit. The most direct route is to run products of
$256\times 256$ and larger on the Arty A7 board, which has 256\,MB of DDR3 memory and an
existing DDR project in the repository. The SoC's external memory port is the natural
connection point, but this has not been tried yet. On the Boolean board, a bounds check should
be added before any larger test, either in the software or in the RTL, so that an oversized
product is rejected rather than silently wrapping around memory.

Performance work comes next. A software baseline is needed for a fair speed-up figure: the
CPU's reference loop should be timed with the RISC-V cycle counter at the same sizes as the
accelerator. Beyond that, the most promising change is a wide, dedicated accelerator port to
block RAM, which would remove the contention with the CPU and raise the memory bandwidth that
currently limits the design. Overlapping loads, computation and stores with ping-pong buffers
would hide much of the remaining latency. Storing the input buffers in distributed RAM instead
of flip-flops would reduce the logic cost of larger \texttt{MAX\_K} values. Finally, the
robustness items should be completed: rejecting configuration writes while a job is running,
checking that addresses stay inside memory, detecting accumulator overflow, aborting cleanly
on a bus error, and adding a completion interrupt.

\end{document}
